% Shared preamble for the Bhagavad Gita PDF (LuaLaTeX)
\documentclass[10pt,a4paper,landscape]{article}
\usepackage[landscape,margin=14mm,top=16mm,bottom=16mm]{geometry}
\usepackage{fontspec}
\usepackage{ltablex}
\usepackage{booktabs}
\usepackage{array}
\usepackage[table]{xcolor}
\usepackage{fancyhdr}
\usepackage{ragged2e}
\usepackage[hidelinks]{hyperref}

\setmainfont{Gentium Plus}
\newfontfamily\devanagarifont{Noto Serif Devanagari}[Script=Devanagari]
\newcommand{\sanskrit}{\devanagarifont}
\newcommand{\iast}{\itshape}

\definecolor{rule}{gray}{0.75}
\definecolor{muted}{gray}{0.4}

\setlength{\parindent}{0pt}
\setlength{\parskip}{4pt}
\keepXColumns
\renewcommand{\arraystretch}{1.25}

\pagestyle{fancy}
\fancyhf{}
\fancyhead[L]{\small\itshape Bhagavad Gita}
\fancyhead[R]{\small\itshape\leftmark}
\fancyfoot[C]{\small\thepage}
\renewcommand{\headrulewidth}{0.3pt}

% \chapter{number}{devanagari name}{transliteration}{english meaning}
\newcommand{\chapter}[4]{%
  \clearpage
  \markboth{Chapter #1 \textperiodcentered{} #3}{}%
  \phantomsection
  \addcontentsline{toc}{section}{Chapter #1: #3 \textperiodcentered{} #4}%
  {\Large\bfseries Chapter #1 \quad {\sanskrit #2} \quad #3}\par
  {\large\color{muted} #4}\par
  \vspace{4pt}%
}
\newcommand{\chaptersummary}[1]{{\small\color{muted}\RaggedRight #1\par}\vspace{6pt}}
\newcommand{\speaker}[1]{{\small\color{muted}#1}}
\newcommand{\verseno}[1]{{\footnotesize\bfseries\color{muted}#1}\newline}
\newcommand{\verserule}{\arrayrulecolor{rule}\hline\arrayrulecolor{black}}

% Header rows of the verse table (tabularx itself is opened in the generated file,
% because it must see a literal \end{tabularx}).
\newcommand{\versehead}{%
  \toprule
  \textbf{Sanskrit} & \textbf{Transliteration} & \textbf{English (\translatorname)} & \textbf{Notes} \\
  \midrule
  \endfirsthead
  \toprule
  \textbf{Sanskrit} & \textbf{Transliteration} & \textbf{English (\translatorname)} & \textbf{Notes} \\
  \midrule
  \endhead
  \bottomrule
  \endlastfoot
}
